\section{Systematik der Gleichrichterbezeichnungen}

Die Bezeichnungen der Gleichrichter folgen einer einheitlichen Nomenklatur:

\begin{itemize}
\item M = Mittelpunktschaltung
\item B = Brückenschaltung
\item Zahl = Pulszahl pro Netzperiode
\item U = ungesteuert (Dioden)
\item C = gesteuert (Thyristoren)
\end{itemize}

\subsection{Bedeutung der Kennzeichnung}

\begin{center}
\begin{tabular}{l|l}
Bezeichnung & Bedeutung \\ \hline
M1U & Einpuls, Mittelpunkt, ungesteuert \\
M1C & Einpuls, Mittelpunkt, gesteuert \\
B2U & Zweipuls, Brücke, ungesteuert \\
B2C & Zweipuls, Brücke, gesteuert \\
B6U & Sechspuls, Brücke, ungesteuert \\
B6C & Sechspuls, Brücke, gesteuert \\
\end{tabular}
\end{center}

\subsection{Grundklassifikation}

\begin{itemize}
\item Ungesteuerte Gleichrichter (U): Dioden, keine Stellmöglichkeit.
\item Gesteuerte Gleichrichter (C): Thyristoren, Stellgrösse Zündwinkel $\alpha$.
\item Pulszahl bestimmt die Restwelligkeit der Gleichspannung.
\end{itemize}

\crd{\textrightarrow\ Diese Systematik gilt für alle folgenden Gleichrichterkapitel.}
